\section{模型的建立与求解：从两束反射到逐次往返}

上一章统一了外角、纵向传播因子与厚度量纲，本章据此从一次返回场出发，将同片双角约束与后续往返逐步写入模型。图~\ref{fig:overall-route} 把这条关系画成同一条数据路径；问题一的两束边界和问题三的往返判据分别落在图\ref{fig:once-return}与图\ref{fig:roundtrip}中。返回阶数随题目需求逐层增加，一次返回先解释条纹，同片双角与完整往返场再分别处理两种材料。%
% src:交接/计划.json:叙事主线
% src:交接/论点脊柱.json:章.10.论点链.0.发现
% src:交接/图注素材_1.json:0.独有信息

三类验证量各有参照：问题一主方法的合成匹配平均误差为 $0.0008\%$，问题二连续留段标准化均方根误差为 $1.21$，问题三同口径误差为 $1.02$。前者衡量已知光学条件下的厚度回收，后两者分别衡量碳化硅与完整往返场的反射率预测表现。%
% src:求解/问题1/结果/合成验证.json:主指标.数值
% src:求解/问题2/结果/验证.json:正式主方法分数_连续留段标准化均方根误差
% src:求解/问题3/结果/回炉轮3候选/仲裁闭环/20260910_131501_81623_建模公式/厚度结果.json:分材料验证.硅.逐折汇总.0.平均标准化均方根误差.两束;分材料验证.硅.逐折汇总.0.平均标准化均方根误差.完整往返;分材料验证.硅.逐折汇总.0.平均标准化均方根误差.训练均值;分材料验证.硅.逐折汇总.0.平均标准化均方根误差.二次趋势;分材料验证.硅.逐折汇总.1.平均标准化均方根误差.两束;分材料验证.硅.逐折汇总.1.平均标准化均方根误差.完整往返;分材料验证.硅.逐折汇总.1.平均标准化均方根误差.训练均值;分材料验证.硅.逐折汇总.1.平均标准化均方根误差.二次趋势;分材料验证.碳化硅.逐折汇总.0.平均标准化均方根误差.两束;分材料验证.碳化硅.逐折汇总.0.平均标准化均方根误差.完整往返;分材料验证.碳化硅.逐折汇总.0.平均标准化均方根误差.训练均值;分材料验证.碳化硅.逐折汇总.0.平均标准化均方根误差.二次趋势;分材料验证.碳化硅.逐折汇总.1.平均标准化均方根误差.两束;分材料验证.碳化硅.逐折汇总.1.平均标准化均方根误差.完整往返;分材料验证.碳化硅.逐折汇总.1.平均标准化均方根误差.训练均值;分材料验证.碳化硅.逐折汇总.1.平均标准化均方根误差.二次趋势

结果关系随模型边界保留：问题二得到的 $10.80\ \mu\mathrm{m}$ 碳化硅拟合厚度，作为问题三继续比较时的碳化硅基准；硅则由完整往返场模型给出 $2.17\ \mu\mathrm{m}$ 的全量厚度估计。碳化硅的修正判断同时核对高阶返回的相干性、空间重叠和仪器分辨率，并比较同片厚度随场模型的变化。%
% src:求解/问题2/结果/厚度结果.json:同片最佳厚度_微米
% src:求解/问题3/结果/回炉轮3候选/仲裁闭环/20260910_131501_81623_建模公式/厚度结果.json:核心指标.碳化硅正式基准厚度_um;核心指标.硅全量完整往返厚度_um

双角原谱高相关曾被用来评价厚度可靠性，核对同片观测关系后，改用双角联动的连续留段检验评估共享参数。正文据此按问题一、问题二、问题三依次展开：先确定两束物理关系，再处理同片双角，最后比较完整往返场及其材料分支。%
% src:交接/实验记录.json:50.现象;50.决定
